% !TEX root = task-2-7-4.tex
% Editable task bank. Append new BankTask blocks at the end of the document.
\documentclass[12pt,a4paper]{article}
\usepackage{physics-textbook}
\usepackage{graphicx,needspace}
\usetikzlibrary{decorations.pathmorphing}
\newif\ifShowSolutions
\ifdefined\TutorEdition
  \ShowSolutionstrue
\else
  \ShowSolutionsfalse % Change to \ShowSolutionstrue to display solutions.
\fi
\newcounter{banktask}
\newsavebox{\BankTaskBox}
\newcommand{\BankTaskContent}[3]{%
  \noindent\textbf{\thebanktask.}\quad #2\par
  \ifShowSolutions
    \Needspace{50mm}% Keep the solution heading with its opening text or diagram.
    #3\par
    \ifdefined\TutorEdition
      \ifcsname TutorSource#1\endcsname
        \nopagebreak\smallskip{\small\csname TutorSource#1\endcsname\par}%
      \fi
    \fi
  \fi
}
\newcommand{\BankTask}[3]{%
  \refstepcounter{banktask}\par\medskip
  \sbox{\BankTaskBox}{\begin{minipage}{\linewidth}%
    \BankTaskContent{#1}{#2}{#3}\end{minipage}}%
  \ifdim\dimexpr\ht\BankTaskBox+\dp\BankTaskBox\relax>\dimexpr\textheight-20mm\relax
    \Needspace{35mm}\BankTaskContent{#1}{#2}{#3}%
  \else
    \Needspace{\dimexpr\ht\BankTaskBox+\dp\BankTaskBox\relax}%
    \noindent\usebox{\BankTaskBox}\par
  \fi
  \medskip
}
\begin{document}
\BookSubsection{2.7.4}{Влажность воздуха}


% Source sdamgia-3218: https://phys-ege.sdamgia.ru/problem?id=3218
\BankTask{sdamgia-3218}{
Давление пара в помещении при температуре $5^\circ\text{C}$ равно $756\,\text{Па}$. Давление насыщенного пара при этой же температуре равно $880\,\text{Па}$. Какова относительная влажность воздуха? (В процентах, округлив до целых.)
}{
\smallskip\noindent\textbf{Решение}\par
Относительная влажность воздуха определяется следующим образом: $\varphi = \frac{p}{p_{\text{н.п.}}} \cdot 100\%,$ где $p$ — давление пара в помещении, а $p_{\text{н.п.}}$ — давление насыщенного пара при той же температуре (эта величина зависит только от температуры воздуха в сосуде). Таким образом, относительная влажность воздуха в помещении равна

$\varphi = \frac{756\text{Па}}{880\text{Па}} \cdot 100\%\approx 86\%.$

Ответ: $86$.
}

% Source sdamgia-9087: https://phys-ege.sdamgia.ru/problem?id=9087
\BankTask{sdamgia-9087}{
Плотность водяного пара в воздухе при температуре $T$ равна $\rho $. Плотность насыщенного водяного пара при этой температуре равна $\rho _{\text{н}}$. Молярная масса воды $\mu _{\text{воды}}=18\,\text{г}/\text{моль}$, а сухого воздуха $\mu_{\text{возд.}}=29\,\text{г}/\text{моль}$.  

\TaskMatchingColumns{Физическая величина}{А)  относительная влажность воздуха при температуре $T$\par Б)  парциальное давление водяного пара при температуре $T$}{Формула}{$1$)  $\frac{\rho R T}{\mu_{\text{воды}}}$\par $2$)  $\frac{\rho_{\text{н}} R T}{\mu_{\text{возд.}}}$\par $3$)  $\rho/ \rho_{\text{н}}$\par $4$)  $\rho_{\text{н}}/ \rho$}\begin{center}\begin{tabular}{|c|c|}\hline А&Б\\\hline &\\\hline\end{tabular}\end{center}
}{
\smallskip\noindent\textbf{Решение}\par
Относительная влажность воздуха связана с парциальным давлением пара при некоторой температуре и давлением насыщенных паров при той же температуре соотношением $\varphi = \frac{p}{p_{\text{н.п.}}} \cdot 100 \%.$ 

Из уравнения Клапейрона — Менделеева давление можно выразить как:

$p= \frac{\nu RT}{V} = \frac{mRT}{V\mu} = \frac{\rho RT}{\mu} .$

Тогда парциальное давление водяного пара будет $p_{\text{вод.п.}}= \frac{\rho RT}{\mu_{\text{воды}}} .$

Из формулы для давления также видно, что соотношение давлений пара и насыщенного пара будут такими же, как и соотношение их плотностей при заданных температурах.

Ответ: $31$.
}

% Source sdamgia-10070: https://phys-ege.sdamgia.ru/problem?id=10070
\BankTask{sdamgia-10070}{
Во вторник и в среду температура воздуха была одинаковой. Парциальное давление водяного пара в атмосфере во вторник было меньше, чем в среду. 

 

$1$. Относительная влажность воздуха во вторник была меньше, чем в среду. 

$2$. Масса водяных паров, содержащихся в $1$ $\text{м}^{3}$ воздуха, во вторник была больше, чем в среду. 

$3$. Плотность водяных паров, содержащихся в воздухе, во вторник и в среду была одинаковой. 

$4$. Давление насыщенных водяных паров во вторник было больше, чем в среду. 

$5$. Концентрация молекул водяного пара в воздухе во вторник была меньше, чем в среду.
}{
\smallskip\noindent\textbf{Решение}\par
Относительная влажность воздуха связана с парциальным давлением пара при некоторой температуре и давлением насыщенных паров при той же температуре соотношением $\varphi = \frac{p}{p_{\text{н.п.}}} \cdot 100 \%.$ 

Проверим правильность утверждений.

$1$. Так как парциальное давление во вторник было ниже, чем в среду, то относительная влажность воздуха во вторник была меньше, чем в среду. Утверждение $1$ — верно.

$2$. Согласно уравнению Менделеева — Клапейрона, при равной температуре и одинаковом объеме большему давлению отвечает большая масса газа. Утверждение $2$ — неверно.

$3$. Так как различается масса водяных паров, содержащихся в воздухе в эти дни, то и плотность водяных паров будет различна во вторник и среду. Утверждение $3$ — неверно.

$4$. При неизменности температуры, давление насыщенного пара остается постоянным. Утверждение $4$ — неверно.

$5$. Концентрация молекул водяного пара в воздухе во вторник была меньше, чем в среду. Утверждение $5$ — верно.

Ответ: $15$.
}

% Source sdamgia-1001: https://phys-ege.sdamgia.ru/problem?id=1001
\BankTask{sdamgia-1001}{
На рисунке представлены два термометра, используемые для определения относительной влажности воздуха с помощью психрометрической таблицы, в которой влажность указана в процентах.

Психрометрическая таблица представлена ниже.


\begin{center}\begin{tikzpicture}[task graph]
\begin{scope}[scale=1.1]\draw(-1,-.3)rectangle(4.5,6);\draw(-1,-.3)--(-1.3,-.7)--(4.8,-.7)--(4.5,-.3);\draw(-1.3,-.7)--(-1.3,-.9)--(4.8,-.9)--(4.8,-.7);\foreach\x/\t in{.3/23,3.1/18}{\draw(\x-1,1.6)rectangle(\x+1,5.6);\node[above]at(\x,5.25){$^\circ\text{C}$};\draw(\x-.1,.6)--(\x-.1,5.2)(\x+.1,.6)--(\x+.1,5.2);\fill[orange!40](\x-.09,.6)rectangle(\x+.09,{2+.1*\t});\foreach\v in{0,1,...,30}{\draw(\x-.3,{2+.1*\v})--(\x+.3,{2+.1*\v});}\foreach\v in{0,5,...,30}{\draw(\x-.5,{2+.1*\v})--(\x+.5,{2+.1*\v});}\foreach\v in{0,10,20,30}{\draw(\x-.75,{2+.1*\v})--(\x+.75,{2+.1*\v});\node[left]at(\x-.45,{2.12+.1*\v}){$\v$};\node[right]at(\x+.45,{2.12+.1*\v}){$\v$};}}\draw[fill=orange!40](.3,.6)circle(.16);\draw(2.4,1.05)..controls(2.6,.85)and(3.6,.85)..(3.8,1.05);\draw(2.4,1.05)..controls(2.6,1.25)and(3.6,1.25)..(3.8,1.05);\fill[orange!30](2.52,.55)--(2.66,-.25)..controls(2.9,-.45)and(3.35,-.45)..(3.55,-.25)--(3.68,.55)--cycle;\draw(2.4,1.05)--(2.66,-.25)..controls(2.9,-.45)and(3.35,-.45)..(3.55,-.25)--(3.8,1.05);\draw(3.1,.55)ellipse(.58 and .08);\draw(3,.6)--(3,.35)--(3.2,.35)--(3.2,.6);\end{scope}
\end{tikzpicture}\end{center}


\begin{center}\small\setlength{\tabcolsep}{4pt}\begin{tabular}{|c|*{9}{c|}}\hline $t_{\text{сух}},\,{}^\circ\text{C}$&\multicolumn{9}{c|}{\shortstack{Разность показаний сухого и влажного\\термометров, ${}^\circ\text{C}$}}\\\hline &0&1&2&3&4&5&6&7&8\\\hline10&100&88&76&65&54&44&34&24&14\\\hline11&100&88&77&66&56&46&36&26&17\\\hline12&100&89&78&68&57&48&38&29&20\\\hline13&100&89&79&69&59&49&40&31&23\\\hline14&100&90&79&70&60&51&42&33&25\\\hline15&100&90&80&71&61&52&44&36&27\\\hline16&100&90&81&71&62&54&45&37&30\\\hline17&100&90&81&72&64&55&47&39&32\\\hline18&100&91&82&73&64&56&48&41&34\\\hline19&100&91&82&74&65&58&50&43&35\\\hline20&100&91&83&74&66&59&51&44&37\\\hline21&100&91&83&75&67&60&52&46&39\\\hline22&100&92&83&76&68&61&54&47&40\\\hline23&100&92&84&76&69&61&55&48&42\\\hline24&100&92&84&77&69&62&56&49&43\\\hline25&100&92&84&77&70&63&57&50&44\\\hline\end{tabular}\end{center}

Какой была относительная влажность воздуха в тот момент, когда проводилась съемка? (В процентах.)
}{
\smallskip\noindent\textbf{Решение}\par
На рисунке видно, что сухой термометр показывает $t_{\text{сух}}=23 ^\circ\text{C},$ а влажный — $t_{\text{вл}}=18 ^\circ\text{C}.$ Следовательно, разность показаний сухого и влажного термометров составляет $t_{\text{сух}} - t_{\text{вл}}=23 ^\circ\text{C} - 18 ^\circ\text{C}=5 ^\circ\text{C}.$ Находим в таблице пересечение соответствующих строки и столбца и получаем, что относительная влажность равна $61\%.$

Ответ: $61$.
}

% Source sdamgia-27946: https://phys-ege.sdamgia.ru/problem?id=27946
\BankTask{sdamgia-27946}{
В жестком закрытом сосуде находится влажный воздух при температуре $16^\circ\text{C}$. Плотность водяных паров в сосуде равна $1{,}078\cdot 10^{-2}\,\text{кг}$/$\text{м}^{3}$. Воздух в сосуде нагревают до $25^\circ\text{C}$. В таблице показана зависимость плотности $\rho_{\text{нп}}$ насыщенных паров воды от температуры $t$. 

\begin{center}\small\setlength{\tabcolsep}{3pt}\begin{tabular}{|c|*{10}{c|}}\hline $t,\,{}^\circ\text{C}$&16&17&18&19&20&21&22&23&24&25\\\hline $\rho_{\text{нп}},\,10^{-2}\,\text{кг}/\text{м}^{3}$&1,36&1,45&1,54&1,63&1,73&1,83&1,94&2,06&2,18&2,30\\\hline\end{tabular}\end{center}

$1$. При температуре $17^\circ\text{C}$ на стенках сосуда есть капли росы.

$2$. Относительная влажность воздуха в сосуде при $18^\circ\text{C}$ равна $70$\%.

$3$. При увеличении температуры относительная влажность воздуха в сосуде уменьшается.

$4$. Давление в сосуде остается постоянным.

$5$. Концентрация водяного пара в сосуде при нагревании остается постоянной.
}{
\smallskip\noindent\textbf{Решение}\par
$1$. Неверно. Температура росы меньше $16^\circ\text{C}$, следовательно, при температуре $17^\circ\text{C}$ роса не выпадет.

$2$. Верно. Относительная влажность находится по формуле: 

$\varphi= \frac{\rho}{\rho_{\text{нп}}} = \frac{1{,}078 \cdot 10 ^{- 2}}{1{,}54 \cdot 10 ^{- 2}} \cdot 100\%=70\%.$

$3$. Верно. При увеличении температуры увеличивается плотность насыщенных водяных паров, плотность водяных паров остается неизменной. Значит, относительная влажность воздуха при увеличении температуры уменьшается.

$4$. Неверно. Из уравнения Клапейрона-Менделеева следует, что $p= \frac{\rho}{M} RT.$ Плотность водяных паров не меняется, а при увеличении температуры давление водяных паров увеличивается.

$5$. Верно. Концентрация водяных паров $n= \frac{N}{V} ,$ число молекул $N=N_A\nu,$ количество вещества $\nu = \frac{m}{M} ,$ масса газа $m=\rho V.$ Откуда $n= \frac{N_A\rho}{M} .$ Значит, при неизменной плотности водяных паров концентрация молекул не меняется.

Ответ: $235$.
}

% Source sdamgia-6819: https://phys-ege.sdamgia.ru/problem?id=6819
\BankTask{sdamgia-6819}{
Один моль влажного воздуха находится в ненасыщенном состоянии при температуре $T$ и давлении $p$. Температуру газа изобарически увеличили. Как изменились при этом относительная влажность воздуха и точка росы при том же давлении? 

$1$. Увеличилась.

$2$. Уменьшилась.

$3$. Не изменилась.

\begin{center}\begin{tabular}{|p{0.43\linewidth}|p{0.43\linewidth}|}\hline
Относительная влажность воздуха & Точка росы\\\hline
\strut & \strut\\\hline\end{tabular}\end{center}
}{
\smallskip\noindent\textbf{Решение}\par
Точка росы — это температура, до которой должен охладиться воздух, чтобы содержащийся в нем пар достиг состояния насыщения и начал конденсироваться в росу. Следовательно, при изобарическом повышении температуры температура точки росы не изменяется.

При увеличении температуры возрастает влагоемкость воздуха, то есть максимальная масса водяного пара, который может содержаться в единице объема воздуха, следовательно, относительная влажность воздуха уменьшается.

Ответ: $23$.
}

% Source sdamgia-8184: https://phys-ege.sdamgia.ru/problem?id=8184
\BankTask{sdamgia-8184}{
Относительная влажность воздуха в закрытом сосуде с поршнем равна $40$\%. Объем сосуда за счет движения поршня медленно уменьшают при постоянной температуре. В конечном состоянии объем сосуда в $3$ раза меньше начального.

$1$. При уменьшении объема сосуда в $2{,}5$ раза на стенках появляется роса.

$2$. Давление пара в сосуде все время увеличивается.

$3$. В конечном и начальном состоянии масса пара в сосуде одинакова.

$4$. При уменьшении объема в $2$ раза относительная влажность воздуха в сосуде стала равна $80$\%.

$5$. В конечном состоянии весь пар в сосуде сконденсировался.
}{
\smallskip\noindent\textbf{Решение}\par
После уменьшения объема в $2$ раза относительная влажность воздуха увеличилась в $2$ раза и стала $80$\%. Когда объем стал в $2{,}5$ раза меньше первоначального, относительная влажность достигла $100$\%, после чего водяные пары начинали конденсироваться на стенках. При дальнейшем уменьшении объема давление водяных паров оставалось постоянным. В конечном состоянии не весь пар в сосуде сконденсировался.

Верны первое и четвертое утверждения.

Ответ: $14$.
}

% Source sdamgia-6911: https://phys-ege.sdamgia.ru/problem?id=6911
\BankTask{sdamgia-6911}{
Два сосуда объемами $20\,\text{л}$ и $30\,\text{л}$, соединенные трубкой с краном, содержат влажный воздух при комнатной температуре. Относительная влажность в сосудах равна соответственно $30$\% и $40$\%. Если кран открыть, то какой будет относительная влажность воздуха в сосудах после установления теплового равновесия, считая температуру постоянной?
}{
\smallskip\noindent\textbf{Решение}\par
Обозначим давление насыщенных водяных паров $p$, тогда парциальные давления паров воды в сосудах: $p_1=\varphi_1 p$ и $p_2 =\varphi_2 p.$ Из уравнения Менделеева — Клапейрона выражаем количество водяных паров в сосудах:

$p_1V_1 = \nu_1 RT ,$ откуда $\nu_1= \frac{p_1V_1}{RT} ,$

$p_2V_2 = \nu_2 RT,$ откуда $\nu_2= \frac{p_2V_2}{RT} .$

После открытия крана парциальное давление водяных паров будет:

$p_0 ( V_1 + V_2 ) = ( \nu_1 + \nu_2 ) RT ,$

$p_0= \frac{p_1V_1 + p_2V_2}{V_1 + V_2} = \frac{\varphi_1 pV_1 + \varphi_2 pV_2}{V_1 + V_2} .$

Тогда относительная влажность в сосудах после установления теплового равновесия:

$\varphi_0= \frac{p_0}{p} = \frac{1}{p} \cdot \frac{\varphi_1 pV_1 + \varphi_2 pV_2}{V_1 + V_2} = \frac{\varphi_1 V_1 + \varphi_2 V_2}{V_1 + V_2} = \frac{30\% \cdot 20 \text{л} + 40\% \cdot 30\text{л}}{50\text{л}} =36\%.$

Ответ: $36$\%.
}

% Source sdamgia-9250: https://phys-ege.sdamgia.ru/problem?id=9250
\BankTask{sdamgia-9250}{
В комнате размером $3\cdot 5$ × $6\,\text{м}$ при температуре $20^\circ\text{C}$ влажность воздуха равна $35$\%. После включения увлажнителя воздуха, производительность которого равна $0{,}36\,\text{л}/\text{ч}$, влажность в комнате стала равна $70$\%. За какое время это произошло? Давление насыщенного пара при $20^\circ\text{C}$ равно $2{,}33\,\text{кПа}$.
}{
\smallskip\noindent\textbf{Решение}\par
Объем воздуха в комнате $V=3 \cdot 5 \cdot 6=90\text{м} ^{3} .$ Масса воды, которая испарена за данное время, равна $m_{\text{в}}=\rho V_{\text{в}}=\rho Pt$

Относительной влажностью называют отношение давления пара к давлению насыщенного пара при той же температуре:

$\varphi = \frac{p}{p_{\text{нп}}} \cdot 100\%.$

Тогда давление водяных паров в комнате в начале и в конце работы увлажнителя принимает значения:

$p_1=\varphi_1\cdot p_{\text{нп}}\quad\text{и}\quad p_2=\varphi_2\cdot p_{\text{нп}}.$

Функция увлажнителя — превратить воду, которая есть в его распоряжении в водяной пар. Тем самым в комнате увеличивается масса водяного пара в воздухе и как следствие увеличивается влажность. Пусть $P$ — производительность увлажнителя, тогда масса водяного пара в комнате до и после связана выражением: 

$m_2=m_1 + P \cdot t \cdot \rho.$

Запишем уравнение Менделеева — Клапейрона для двух состояний водяного пара:

$\begin{cases}p_1V=\frac{m_1}{M}RT,\\p_2V=\frac{m_2}{M}RT.\end{cases}$

Вычтем из нижнего уравнения верхнее:

$p_2V - p_1V= \frac{RT}{M} ( m_2 - m_1 ) = \frac{RT}{M} \cdot P \cdot t \cdot \rho.$

Выразим отсюда время:

$t= \frac{VM ( p_2 - p_1 )}{RT \cdot P \cdot \rho} = \frac{VMp_{\text{нп}} ( \varphi_2 - \varphi_1 )}{RT \cdot P \cdot \rho} = \frac{90 \cdot 18 \cdot 10 ^{- 3} \cdot 2{,}33 \cdot 10 ^{3} ( 0{,}70 - 0{,}35 )}{8{,}31 \cdot 293 \cdot 0{,}36 \cdot 10 ^{- 3} \cdot 1000} \approx 1{,}5\text{ч}.$

Ответ: $1{,}5$ ч.
}

% Source sdamgia-32562: https://phys-ege.sdamgia.ru/problem?id=32562
\BankTask{sdamgia-32562}{
В закрытом сосуде объемом V  =  $30$ л находится влажный воздух массой m  =  $45\,\text{г}$ при температуре t  =  $80$ °C и давлении p  =  $160$ кПа. Определите относительную влажность воздуха, находящегося в сосуде. Давление $p_{\,\text{н}}$ насыщенных паров воды при температуре $80$ °С равно $47$ кПа.
}{
\smallskip\noindent\textbf{Решение}\par
Относительная влажность воздуха определяется по формуле $\varphi = \frac{p_{\,\text{п}}}{p_{\,\text{н}}} \cdot 100\%.$
По закону Дальтона давление в сосуде складывается из парциального давления водяного пара $p_{\,\text{п}}$ и давления сухого воздуха $p_{\,\text{возд}}$, то есть $p=p_{\,\text{п}} + p_{\,\text{возд}}.$
Из уравнения Клапейрона  — Менделеева найдем давление сухого воздуха:

$p_{\,\text{возд}}V= \frac{m_{\,\text{возд}}RT}{M_{\,\text{возд}}} ,$ следовательно,  $p_{\,\text{возд}}= \frac{m_{\,\text{возд}}RT}{VM_{\,\text{воз}}} .$

Аналогично находим парциальное давление водяного пара:

$p_{\,\text{п}}V= \frac{m_{\,\text{п}}RT}{M_{\,\text{вод}}} ,$ следовательно, $p_{\,\text{п}}= \frac{m_{\,\text{п}}RT}{VM_{\,\text{вод}}} .$

Учтем, что масса влажного воздуха равна $m=m_{\,\text{п}} + m_{\,\text{возд}}.$ Объединяя уравнения, выразим парциальное давление водяного пара:

\[p_{\,\text{п}}= \frac{pVM_{\,\text{возд}} - mRT}{V ( M_{\,\text{возд}} - M_{\,\text{вод}} )} .\]

Тогда относительная влажность воздуха равна:

\[\varphi = \frac{pVM_{\,\text{возд}} - mRT}{p_{\,\text{н}}V ( M_{\,\text{возд}} - M_{\,\text{вод}} )} \cdot 100\%,\]

 то есть:

\[\varphi = \frac{160 \cdot 10 ^3 \cdot 30 \cdot 10 ^{ - 3 } \cdot 29 \cdot 10 ^{ - 3 } - 45 \cdot 10 ^{ - 3 } \cdot 8{,}31 \cdot 353}{47 \cdot 10 ^3 \cdot 30 \cdot 10 ^{ - 3 } \cdot ( 29 \cdot 10 ^{ - 3 } - 18 \cdot 10 ^{ - 3 } )} \cdot 100\%\approx 46\%.\]

Ответ: $46$\%.
}


% Source 430072: https://ege.fipi.ru/bank/questions.php?search=1&proj=BA1F39653304A5B041B656915DC36B38&qid=430072
\BankTask{430072}{
Относительная влажность воздуха в закрытом сосуде с поршнем равна \(40\,\%\). Объём сосуда за счёт движения поршня медленно уменьшают при постоянной температуре. В конечном состоянии объём сосуда в \(4\) раза меньше начального.
\begin{enumerate}
\item При уменьшении объёма сосуда в \(2{,}5\) раза на стенках появляется роса.
\item Давление пара в сосуде всё время увеличивается.
\item В конечном и начальном состояниях масса пара в сосуде одинакова.
\item При уменьшении объёма в \(2\) раза относительная влажность воздуха в сосуде стала равна \(80\,\%\).
\item В конечном состоянии весь пар в сосуде сконденсировался.
\end{enumerate}
}{
\smallskip\noindent\textbf{Решение от GPT}\par
При постоянной температуре давление насыщенного пара \(p_{\text{нас}}\) постоянно. До насыщения \(pV=\mathrm{const}\), поэтому \(\varphi=0{,}4V_0/V\).
\begin{enumerate}
\item Верно: насыщение достигается при \(V=0{,}4V_0=V_0/2{,}5\); начинается выпадение росы.
\item Неверно: после насыщения давление пара остаётся равным \(p_{\text{нас}}\).
\item Неверно: часть пара конденсируется. При \(V_f=V_0/4\) отношение масс пара равно \(m_f/m_0=(1/4)/0{,}4=0{,}625\).
\item Верно: при \(V=V_0/2\) влажность \(\varphi=0{,}8=80\,\%\).
\item Неверно: в конечном состоянии остаётся насыщенный пар.
\end{enumerate}
Ответ: \(14\).
}

% Source DE8B36: https://ege.fipi.ru/bank/questions.php?search=1&proj=BA1F39653304A5B041B656915DC36B38&qid=DE8B36
\BankTask{DE8B36}{
В понедельник и вторник температура воздуха была одинаковой. Парциальное давление водяного пара в атмосфере в понедельник было меньше, чем во вторник.
\begin{enumerate}
\item Масса водяных паров, содержащихся с $1\ \text{м}^{3}$ воздуха в понедельник была больше, чем во вторник
\item Относительная влажность воздуха в понедельник была меньше, чем во вторник
\item Концентрация молекул водяного пара в воздухе в понедельник и вторник была одинаковой
\item Давление насыщенных водяных паров в понедельник было больше, чем во вторник
\item Плотность водяных паров, содержащихся в воздухе в понедельник была меньше, чем во вторник
\end{enumerate}
}{
\smallskip\noindent\textbf{Решение от GPT}\par
При одинаковой температуре $p_{\text{нас}}(T)$ одинаково в оба дня. По формулам $\varphi=p/p_{\text{нас}}$, $n=p/(kT)$ и $\rho=Mp/(RT)$ меньшему парциальному давлению соответствуют меньшие влажность, концентрация и плотность пара.
\begin{enumerate}
\item Неверно: масса пара в $1\ \text{м}^{3}$ в понедельник меньше.
\item Верно: $\varphi_{\text{пн}}<\varphi_{\text{вт}}$.
\item Неверно: $n_{\text{пн}}<n_{\text{вт}}$.
\item Неверно: давления насыщенного пара одинаковы.
\item Верно: $\rho_{\text{пн}}<\rho_{\text{вт}}$.
\end{enumerate}
Ответ: $25$.
}

% Source 055098: https://ege.fipi.ru/bank/questions.php?search=1&proj=BA1F39653304A5B041B656915DC36B38&qid=055098
\BankTask{055098}{
Стальную кастрюлю, наполовину заполненную водой, поставили на плиту, закрыли лёгкой крышкой и стали нагревать воду до кипения. Атмосферное давление составляет \(760\,\text{мм рт. ст.}\)
\begin{enumerate}
\item Относительная влажность воздуха под крышкой в процессе нагревания воды уменьшается.
\item Внутренняя энергия водяного пара в кастрюле в процессе нагревания воды увеличивается.
\item При нагревании воды выше \(80\,^{\circ}\mathrm{C}\), когда схлопываются образующиеся пузырьки, давление водяных паров постепенно уменьшается.
\item Температура дна кастрюли вблизи границы с водой при кипении немного превышает \(100\,^{\circ}\mathrm{C}\).
\item Плотность насыщенных водяных паров над поверхностью воды в процессе её нагревания до температуры кипения увеличивается.
\end{enumerate}
}{
\smallskip\noindent\textbf{Решение от GPT}\par
1) Неверно: под крышкой воздух насыщается водяным паром, его относительная влажность не уменьшается.

2) Верно: при нагревании растут температура и количество водяного пара, поэтому его внутренняя энергия увеличивается.

3) Неверно: давление насыщенного пара с повышением температуры увеличивается; пузырьки схлопываются при попадании в более холодные слои воды вследствие конденсации пара.

4) Верно: при данном атмосферном давлении вода кипит при \(100\,^{\circ}\mathrm{C}\); для передачи теплоты от дна воде температура дна должна быть немного выше.

5) Верно: плотность насыщенного водяного пара увеличивается с температурой.

Ответ: \(245\).
}

% Source 6ECE95: https://ege.fipi.ru/bank/questions.php?search=1&proj=BA1F39653304A5B041B656915DC36B38&qid=6ECE95
\BankTask{6ECE95}{
В жёстком герметичном сосуде объёмом \(1\,\text{м}^{3}\) при температуре \(289\,\text{К}\) длительное время находился влажный воздух и \(10\,\text{г}\) воды. Сосуд медленно нагрели до температуры \(298\,\text{К}\).

Таблица плотности насыщенных паров воды:
\begin{center}
\begin{tabular}{|c|c|c|c|c|c|c|c|c|c|c|}
\hline
$t,\,{}^\circ\text{C}$ & $16$ & $17$ & $18$ & $19$ & $20$ & $21$ & $22$ & $23$ & $24$ & $25$ \\
\hline
$\rho_{\text{нп}},\,10^{-2}\,\text{кг}/\text{м}^3$ & $1{,}36$ & $1{,}45$ & $1{,}54$ & $1{,}63$ & $1{,}73$ & $1{,}83$ & $1{,}94$ & $2{,}06$ & $2{,}18$ & $2{,}30$ \\
\hline
\end{tabular}
\end{center}
\begin{enumerate}
\item При температуре \(23\,^{\circ}\mathrm{C}\) влажность воздуха в сосуде была равна \(48{,}5\,\%\).
\item В течение всего опыта в сосуде находилась вода в жидком состоянии.
\item Так как объём сосуда не изменялся, давление влажного воздуха увеличивалось пропорционально его температуре.
\item В начальном состоянии при температуре \(289\,\text{К}\) пар в сосуде был насыщенный.
\item Парциальное давление сухого воздуха в сосуде не изменялось.
\end{enumerate}
}{
\smallskip\noindent\textbf{Решение от GPT}\par
Начальная температура \(t_1=16\,^{\circ}\mathrm{C}\), конечная \(t_2=25\,^{\circ}\mathrm{C}\). Вначале пар насыщенный, поскольку длительное время находился в равновесии с жидкой водой. Его масса \(m_{\text{п}1}=\rho_{\text{нп}1}V=13{,}6\,\text{г}\). Общая масса воды и пара равна \(23{,}6\,\text{г}\). При \(25\,^{\circ}\mathrm{C}\) для насыщения требуется \(23{,}0\,\text{г}\) пара, поэтому остаётся \(0{,}6\,\text{г}\) жидкости.

1) Неверно: при \(23\,^{\circ}\mathrm{C}\) жидкость ещё есть, пар насыщенный, \(\varphi=100\,\%\).

2) Верно: жидкая вода сохраняется даже при конечной температуре.

3) Неверно: при испарении количество вещества газовой смеси увеличивается, поэтому её давление не пропорционально абсолютной температуре.

4) Верно: начальный пар насыщенный.

5) Неверно: \(p_{\text{сух}}=\nu_{\text{сух}}RT/V\), поэтому \(p_{\text{сух}2}/p_{\text{сух}1}=298/289>1\).

Ответ: \(24\).
}

% Source 25411: https://phys-ege.sdamgia.ru/problem?id=25411
\BankTask{25411}{
Сосуд разделен тонкой перегородкой на две части, отношение объемов которых $\frac{V_2}{V_1} =3.$ В первой части сосуда находится воздух с относительной влажностью $\varphi_1=80\%.$ Если убрать перегородку, то относительная влажность воздуха в сосудах после установления теплового равновесия будет равна $50$\%. Определите начальную влажность во втором сосуде, считая, что температура воздуха в частях сосуда одинакова и не изменилась после снятия перегородки.
}{
\smallskip\noindent\textbf{Решение}\par
Найдем давление в каждой части сосуда из формулы относительной влажности воздуха и давление в сосуде после открытия крана:$p_1=\varphi_1 p_{\,\text{нп}},$ $p_2=\varphi_2 p_{\,\text{нп}},$ $p=\varphi p_{\,\text{нп}}.$
Запишем закон сохранения энергии $U_1 + U_2=U,$ где внутренняя энергия определяется по формуле $U= \frac{i}{2} \nu RT.$ Учитывая, что $p_1V_1=\nu_1 RT,$ $p_2V_2=\nu_2 RT,$ получаем:

$\varphi_1 p_{\,\text{нп}}V_1 + \varphi_2 p_{\,\text{нп}}V_2=\varphi p_{\,\text{нп}} ( V_1 + V_2 ) ,$ откуда $\varphi_1 V_1 + 3\varphi_2 V_1=4\varphi V_1;$

$\varphi_2= \frac{4\varphi - \varphi_1}{3} = \frac{4 \cdot 50 - 80}{3} =40\%.$

Ответ: $40$\%.
}

% Source 3655: https://phys-ege.sdamgia.ru/problem?id=3655
\BankTask{3655}{
В закрытом цилиндрическом сосуде находится влажный воздух при температуре $100$ °C. Для того, чтобы на стенках этого сосуда выпала роса, требуется изотермически изменить объем сосуда в $25$ раз. Чему приблизительно равна первоначальная абсолютная влажность воздуха в сосуде? Ответ приведите в граммах на кубический метр, округлите до целых.
}{
\smallskip\noindent\textbf{Решение}\par
Абсолютная влажность воздуха  — это физическая величина, показывающая массу водяных паров, содержащихся в $1$ $\,\text{м} ^3$ воздуха. Другими словами, это плотность водяного пара в воздухе. На стенках сосуда при изотермическом сжатии начнет образовываться роса после того, как пар достигнет состояния насыщения. Как известно, кипение начинается, когда давление насыщенных паров сравнивается с внешним давлением. Таким образом, давление насыщенных паров при $100$ °C равно нормальному атмосферному давлению, то есть порядка $p_{\,\text{н}}. \,\text{п}.=10 ^{5} \,\text{Па}.$
Определим сперва, какое давление имеет пар до начала сжатия. Водяной пар подчиняется уравнениям идеального газа, в частности, при изотермическом сжатии выполняется закон Бойля  — Мариотта: $pV=const.$ Следовательно, начальное давление водяного пара было в $25$ раз меньше $p_{\,\text{н}}. \,\text{п}.=10 ^{5} \,\text{Па}$ и равнялось $p_0=4000\,\text{Па}.$
Определим теперь первоначальную абсолютную влажность. Для этого воспользуемся уравнением состояния Клапейрона  — Менделеева $pV= \frac{m}{\mu} RT .$ Откуда абсолютная влажность воздуха была равна

\[\rho= \frac{m}{V} = \frac{p\mu}{RT} = \frac{4000\,\text{Па} \cdot 0{,}018\,\text{кг/моль}}{8{,}31\,\text{Дж}/ ( \,\text{моль} \cdot \,\text{К} ) \cdot 373\,\text{К}} \approx 0{,}023 \,\text{кг/м} ^3 =23\,\text{г/м} ^3 .\]

Ответ: $23\,\text{г}$/м$^{3}$.
}

% Source 19871: https://phys-ege.sdamgia.ru/problem?id=19871
\BankTask{19871}{
В комнате при $20$ °С относительная влажность $40$\%. При физической нагрузке через легкие человека приходит $15$ л воздуха за $1$ минуту. Выдыхаемый воздух имеет температуру $34$ °C и влажность $100$\%. Давление насыщенного пара при $20$ °С равно $2{,}34$ кПа, а при $34$ °C составляет $5{,}32$ кПа. Какую массу воды теряет тело человека за $1$ час за счет дыхания? Считать, что выдыхаемый воздух имеет такой же объем, какой проходит через легкие человека. Влажность воздуха в комнате не меняется.
}{
\smallskip\noindent\textbf{Решение}\par
$1$.  Парциальное давление пара при температуре t  =  $20$ °C равно $p=\varphi p_{\,\text{нп}1}.$ Из уравнения Клапейрона  — Менделеева $pV= \frac{m_1}{M} RT_1$ находим массу пара в первом состоянии $m_1= \frac{\varphi p_{\,\text{нп}1}VM}{RT_1} .$ Аналогично для пара во втором состоянии $m_2= \frac{\varphi p_{\,\text{нп}2}VM}{RT_2} ,$ его масса равна $m_2= \frac{p_{\,\text{нп}2}VM}{RT_2} .$
За время $\tau=1\,\text{ч} =60 \,\text{мин}$ человек выдыхает объем водяного пара, равный $V= \frac{\Delta V}{\Delta\tau} \cdot \tau,$ где по условию задачи $\frac{\Delta V}{\Delta\tau} =15 \cdot 10 ^{ - 3 } \,\text{м} ^3 /\,\text{мин}.$ Вычислим объем: V  =  $15$ · $10^{-3}$ · $60$  =  $0{,}9\,\text{м}^{3}$.
$3$.  За $1$ час масса потерянной человеком воды равна:

\[\Delta m=m_2 - m_1= \frac{VM}{R} ( \frac{p_{\,\text{нп}2}}{T_2} - \frac{\varphi p_{\,\text{нп}1}}{T_1} ) = \frac{VM ( p_{\,\text{нп}2}T_1 - \varphi p_{\,\text{нп}1}T_2 )}{RT_1T_2} ,\]

\[\Delta m= \frac{0{,}9 \cdot 18 \cdot 10 ^{ - 3 } ( 5{,}32 \cdot 293 - 2{,}34 \cdot 0{,}4 \cdot 307 ) \cdot 10 ^3}{8{,}31 \cdot 307 \cdot 293} \approx 0{,}028\,\text{кг}=28\,\text{г}.\]

Ответ: $28\,\text{г}$.
}

% Source 38662: https://phys-ege.sdamgia.ru/problem?id=38662
\BankTask{38662}{
Вертикальный цилиндрический сосуд с влажным воздухом влажностью $\varphi$ закрыт невесомым поршнем площадью S, способным скользить без трения. В начальном состоянии поршень находится в равновесии. На поршень медленно насыпают песок. Определите минимальную массу песка m, при которой на стенках сосуда появится роса. Температура в сосуде поддерживается постоянной. Внешнее атмосферное давление равняется $p_{0}$.
}{
\smallskip\noindent\textbf{Решение}\par
В первом состоянии поршень невесомый, поэтому давление влажного воздуха равно атмосферному давлению $p_1=p_0.$ Из формулы влажности $\varphi = \frac{p_{\,\text{п}1}}{p_{\,\text{нп}}} ,$ где $p_{\,\text{п}1}$  — парциальное давление водяного пара в первом состоянии, $p_{\,\text{нп}}$  — давление насыщенного пара при данной температуре.
Во втором состоянии выпала роса, что означает  — пар стал насыщенным. Следовательно,  в сосуде давление водяных паров $p_{\,\text{п}2}$ стало равным $p_{\,\text{нп}}$. 
Поршень в равновесии, следовательно, $p_2 = p_0 + \frac{mg}{S} ,$  где $p_{2}$  — давление пара после насыпания песка.
И для ненасыщенного водяного пара, и для влажного воздуха можно применить закон Бойля-Мариотта при изотермическом процессе:

$\frac{V_2}{V_1} = \frac{p_1}{p_2}$
$\frac{V_2}{V_1} = \frac{p_{\,\text{п}1}}{p_{\,\text{п}2}} ,$

где $V_{1}$ и $V_{2}$ начальный и конечные объемы влажного воздуха в сосуде под поршнем. Из этих соотношений получаем:

\[\frac{p_1}{p_2} = \frac{p_{\,\text{п}1}}{p_{\,\text{п}2}} .\]

 Объединяя эти уравнения, получаем, что масса песка равна

\[m= \frac{p_0S ( 1 - \varphi )}{\varphi g} .\]

Ответ: $m = \frac{p_0 S}{g} \cdot ( \frac{1 - \varphi}{\varphi} ) .$
}

% Source 30001: https://phys-ege.sdamgia.ru/problem?id=30001
\BankTask{30001}{
В запаянной с одного конца трубке находится влажный воздух, отделенный от атмосферы столбиком ртути длиной l  =  $76\,\text{мм}$. Когда трубка лежит горизонтально, относительная влажность воздуха $\varphi_1$ в ней равна $80$\%. Какой станет относительная влажность этого воздуха $\varphi_2,$ если трубку поставить вертикально, открытым концом вниз? Атмосферное давление равно $760\,\text{мм}$ рт. ст. Температуру считать постоянной.
}{
\smallskip\noindent\textbf{Решение}\par
$1$.  Относительная влажность воздуха в трубке в первом случае: $\varphi_1= \frac{p_1}{p_{\,\text{н}}} ,$ где $p_{1}$  — давление водяных паров в трубке, $p_{\,\text{н}}$  — давление насыщенных водяных паров при той же температуре. Для второго случая относительная влажность воздуха в трубке: $\varphi_2= \frac{p_2}{p_{\,\text{н}}} ,$ где $p_{2}$  — парциальное давление водяных паров в трубке.

\begin{center}\begin{tikzpicture}[task graph,>=stealth]\draw(.35,0)arc(-90:-270:.35)--(4.4,.7);\draw(.35,0)--(4.4,0);\draw(4.4,.35)ellipse(.16 and .43);\draw[fill=orange!30](1.62,0)rectangle(3.05,.7);\draw(0,-.1)--(0,-1.05)(1.62,0)--(1.62,-1.05)(3.05,0)--(3.05,-1.05);\draw[<->](0,-.88)--(1.62,-.88)node[midway,above]{\(l_1\)};\draw[<->](1.62,-.88)--(3.05,-.88)node[midway,above]{\(l\)};\node at(3.96,.95){\(p_0\)};\end{tikzpicture}\end{center}

\begin{center}\begin{tikzpicture}[task graph,>=stealth]\draw(0,0)--(0,3.55)arc(180:0:.33)--(.66,0);\draw(.33,0)ellipse(.44 and .15);\draw[fill=orange!30](0,1.26)rectangle(.66,2.48);\draw(.66,1.26)--(1.75,1.26)(.66,2.48)--(1.75,2.48)(.33,3.88)--(1.75,3.88);\draw[<->](1.59,1.26)--(1.59,2.48)node[midway,left]{\(l\)};\draw[<->](1.59,2.48)--(1.59,3.88)node[midway,left]{\(l_2\)};\node at(-.46,.4){\(p_0\)};\end{tikzpicture}\end{center}

$2$.  Давление влажного воздуха в первом случае равно атмосферному давлению: $p_1 \,\text{вл}=p_0,$ а во втором случае  — $p_2\,\text{вл}=p_0 - \rho g l,$ где \ensuremath{\rho} — плотность ртути, а l  — длина столбика ртути. Заметим, что $p_0=\rho g H,$ где H  =  $760\,\text{мм}$. 
$3$.  Объем влажного воздуха при изменении положения трубки изменился так, как изменились длины столбиков воздуха, запертых ртутью. Считаем влажный воздух идеальным газом. Поэтому по закону Бойля  — Мариотта можем записать:

\[\frac{V_2}{V_1} = \frac{l_2}{l_1} = \frac{p_0}{p_0 - \rho g l} .\]

$4$.  Из уравнения Клапейрона  — Менделеева следует, что: 

$p= \frac{\nu R T}{V} ,$ $p_{\,\text{вл}}= \frac{( \nu + \nu_{\,\text{св}} ) R T}{V} ,$

 где p  — парциальное давление водяного пара, $p_{\,\text{вл}}$  — давление влажного воздуха, $\nu$ — количество моль водяного пара в трубке, $\nu_{\,\text{св}}$  — количество моль сухого воздуха в трубке. Отсюда следует, что:

\[\frac{p}{p_{\,\text{вл}}} = \frac{\nu}{\nu + \nu_{\,\text{св}}} = const.\]

Поэтому $\frac{p_1}{p_2} = \frac{p_1 \,\text{вл}}{p_2 \,\text{вл}} ,$ и для изотермического процесса для парциальных давлений водяного пара имеем: 

\[\frac{p_1}{p_2} = \frac{V_2}{V_1} = \frac{p_0}{p_0 - \rho g l} .\]

$5$.  После объединения записанных выше уравнений получим выражение для относительной влажности воздуха во втором случае: 

\[\varphi_2=\varphi_1 \cdot \frac{p_0 - \rho g l}{p_0} =\varphi_1 \cdot \frac{\rho g H - \rho g l}{\rho g H} =\varphi_1 \frac{H - l}{H} =80 \cdot \frac{760 - 76}{760} =72 \%.\]

Ответ: $\varphi_2=72 \%.$
}

% Source 50825: https://phys-ege.sdamgia.ru/problem?id=50825
\BankTask{50825}{
Сосуд объёмом V  =  $2$ л разделён на две части. В одной части объёмом $V_{1}$ находятся вода и насыщенный водяной пар. Масса воды в этой части в $2$ раза меньше массы пара. После удаления перегородки пар заполняет весь сосуд. Температура во время процесса остаётся постоянной. После установления равновесия относительная влажность в сосуде составляет $30$\%. Определите объём $V_{2}$.
}{
\smallskip\noindent\textbf{Решение}\par
Отметим, что в сосуде первоначально находился насыщенный пар, следовательно, относительная влажность воздуха была $100$\%. Так как температура во время процесса не изменялась, то плотность насыщенного пара оставалась неизменной и равной первоначальной.
В сосуде объемом $V_{1}$ находились вода и пар. Следовательно, общая масса содержимого равнялась 

\[m=m_{\,\text{п}} + m_{\,\text{в}}=m_{\,\text{п}} + \frac{m_{\,\text{п}}}{2} = \frac{3}{2} m_{\,\text{п}}.\]

Плотность пара была равна плотности насыщенного пара $\rho_{\,\text{н}}.$ Учитывая, что масса пара $m_{\,\text{п}}=\rho_{\,\text{н}}V_1,$ получаем, что масса содержимого в первом сосуде равнялась $m= \frac{3}{2} \rho_{\,\text{н}}V_1.$
После удаления перегородки масса содержимого не изменилась. Так как пар стал ненасыщенным, то вся вода испарилась, в сосуде был пар с плотностью $\rho,$ которая стала равной $\rho= \varphi \rho_{\,\text{н}}.$ При этом сохранившаяся масса пара равна $m=\rho V= \varphi \rho_{\,\text{н}}V.$
Учитывая неизменность массы, получаем равенство $\frac{3}{2} \rho_{\,\text{н}}V_1=\varphi \rho_{\,\text{н}}V,$ откуда находим объем первой части сосуда:

\[V_1= \frac{2 \varphi V}{3} .\]

Следовательно, объем второй части сосуда

\[V_2=V - V_1=V - \frac{2 \varphi V}{3} = \frac{3 - 2 \varphi}{3} V= \frac{3 - 2 \cdot 0{,}3}{3} \cdot 2 =1{,}6\,\text{л}.\]

Ответ: $1{,}6$ л.
}

% Source 25414: https://phys-ege.sdamgia.ru/problem?id=25414
\BankTask{25414}{
В закрытом сосуде при температуре $100 ^\circ$ находится влажный воздух с относительной влажностью воздуха $60$\% под давлением $100$ кПа. Объем сосуда изотермически уменьшили в $2{,}5$ раза. Во сколько раз нужно увеличить абсолютную температуру воздуха без изменения объема сосуда, чтобы получить такое же конечное давление? Объемом сконденсировавшейся воды пренебречь.
}{
\smallskip\noindent\textbf{Решение}\par
Влажный воздух представляет собой смесь воздуха и водяных паров. Так как температура воздуха $100 ^\circ\,\text{С} ,$ то давление насыщенного водяного пара равно $100$ кПа, в данном случае  — давлению влажного воздуха р.
По закону Дальтона давление газа в сосуде равно сумме парциальных давлений: $p=p_1\,\text{возд} + p_{\,\text{пара}}.$
Из формулы относительной влажности воздуха следует, что давление пара $p_{\,\text{пара}}=\varphi p_{\,\text{н}}. \,\text{п}..$ Откуда давление сухого воздуха в сосуде $p_1\,\text{возд}=p ( 1 - \varphi ) .$
После изотермического сжатия водяной пар стал насыщенным, его давление стало равным давлению насыщенного пара. Давление сухого воздуха при изотермическом сжатии в $2{,}5$ раза увеличилось в $2{,}5$ раза по закону Бойля-Мариотта $pV=const,$ $p_2\,\text{возд}=2{,}5p_1\,\text{возд}.$ Тогда давление влажного воздуха в сосуде стало $p_2=p_2\,\text{возд} + p_{\,\text{н}}. \,\text{п}..$ Таким образом, $p_2=2{,}5 ( 1 - \varphi ) p + p= ( 3{,}5 - 2{,}5\varphi ) p.$
При изохорном нагревании по закону Шарля $\frac{p}{T} =const.$ Откуда:

\[\frac{T_2}{T_1} = \frac{p_2}{p} =3{,}5 - 2{,}5\varphi=3{,}5 - 2{,}5 \cdot 0{,}6=2.\]

Значит, абсолютную температуру воздуха нужно увеличить в $2$ раза.

Ответ: в $2$ раза.
}

\end{document}
